\paragraph{Дислокация над включением.} Пластина того же сечения,
$126{,}0\times124{,}0\times125{,}5$~\AA{} (117\,363 атома), содержит сферу
Al$_{13}$Fe$_4$ радиусом 25~\AA{} в нижней половине; плоскость скольжения
проходит на 12~\AA{} выше верхней точки сферы, и при построении центр сферы
стоит на 15~\AA{} впереди линии вдоль $x$. Включение закреплено, как в
расчётах поля, верхний и нижний блоки жёсткие, как в расчёте без включения.
Энергия матрицы минимизируется методом сопряжённых градиентов до нормы сил
$10^{-3}$~эВ/\AA{} без прогрева, затем верхний блок
сдвигается по $z$, пока среднее $\sigma_{zz}$ матрицы не станет близким
к нулю. При этой релаксации линия отходит от включения на «DISL_RELAX»~\AA{}
и останавливается в «DISL_D»~\AA{} от центра сферы вдоль $x$. Положение линии
равно среднему положению двух частичных дислокаций, найденных методом DXA;
расстояние между ними во всех состояниях «DISL_SPLIT»~\AA. Затем атомы включения
получают деформацию $+\varepsilon$, и матрица релаксирует заново: линия
отходит от включения ещё на «DISL_SHIFT»~\AA. После возвращения атомов
включения в прежние положения и новой релаксации линия стоит в «DISL_REV»~\AA{}
от исходного положения.

Нагружение этих пластин сдвигом не даёт отдельной оценки влияния поля.
Уже без нагрузки верхний и нижний блоки несут касательные напряжения
«BLOCK_TOP» и «BLOCK_BOT»~МПа; разницу передаёт закреплённое включение.
Поэтому напряжение, при котором линия срывается с места, зависит прежде всего
от закрепления включения, и вклад его деформации из этих расчётов выделить нельзя.
